
%!TEX TS-program = xelatex

\documentclass[letterpaper,oneside,11pt,article, portrait]{memoir}
\usepackage[pdftitle={Scientific Paper Checklist}, pdfauthor={Trevor Rife}, colorlinks=true, urlcolor=blue]{hyperref}
\usepackage{colortbl}
\usepackage{color}
\usepackage{multirow}

\setlrmarginsandblock{.8in}{.8in}{*}
\setulmarginsandblock{.8in}{.8in}{*}

\usepackage{array,ragged2e}
\usepackage{fontspec,xunicode}
\defaultfontfeatures{Mapping=tex-text}


\definecolor{gray}{rgb}{0.6,0.6,0.6}
\newcommand{\doi}[1]{doi:\href{http://dx.doi.org/#1}{#1}}

\newcommand{\journal}[1]{\textit{#1}} 			% journals

\makepagestyle{footer}
\makeevenfoot{footer}{\footnotesize{Rife Lab Paper Checklist}}{}{\thepage} 
\makeoddfoot{footer}{\footnotesize{Rife Lab Paper Checklist}}{}{\thepage}

\setsecnumdepth{section}

\checkandfixthelayout	% for memoir class

\begin{document}
\pagestyle{empty}


\begin{table}[t]

\centering

{\Large Checklist for Scientific Papers (work in progress)}

 \setlength{\extrarowheight}{.2in}
\begin{tabular}{| p{3.5in} | c | c | c | l |}


\hline
\textbf{Item}& \textbf{Yes}& \textbf{No}& \textbf{Why not?*} & \textbf{See section} \\ \hline

% Open
\rowcolor[gray]{.9} \multicolumn{5}{| l |}{Openness} \\ \hline

Made data publicly available?	& & & & \multirow{2}{*}{\ref{open} (page \pageref{open})}\\ \cline{1-4}

Made data analyses publicly available? & & & &  \\  \hline

Submitted a preprint?& & & & \ref{preprint} (page \pageref{preprint}) \\ \hline

Published in an open access journal? & & & & \ref{openaccess} (page \pageref{openaccess}) \\ \hline

Published original figures \href{http://creativecommons.org/licenses/by/4.0/}{CC-by} prior to journal? & & & & \ref{CCFigure} (page \pageref{CCFigure}) \\ \hline

% Statistics
\rowcolor[gray]{.9} \multicolumn{5}{| l |}{Statistics} \\ \hline

Explained all conducted tests? I.e., included a statement such as ``We report how we determined our sample size, all data exclusions (if any), all manipulations, and all measures in the study.''	& & & & \ref{explainAllTests} (page \pageref{explainAllTests})  \\ \hline

Collected an informative data set? & & & & \ref{n} (page \pageref{n}) \\ \hline

Plotted the data informatively? & & & & \ref{plotting} (page \pageref{plotting}) \\ \hline

Avoided p values?	& & & & \ref{avoidNHST} (page \pageref{avoidNHST})  \\ \hline

Included effect sizes? 	& & & & \ref{avoidNHST} (page \pageref{avoidNHST})  \\ \hline

Included confidence intervals? & & & & \ref{CI} (page \pageref{CI}) \\ \hline

\end{tabular}

{\scriptsize * D = don't know how, H = too hard, S = scary, E = ethical constraints make it impossible, B = bad for science, O = other (specify)}

\end{table}

\clearpage
%\newpage

\pagestyle{footer}

\noindent {\itshape The goal of this checklist is to promote discussion about---and implementation of---best practices in research. This checklist is built from a similar document created by \href{http://github.com/jpeelle/paperchecklist}{Jonathan Peelle}}.

%%%%%%%%%%%%%%%%%%%%%%%%
\chapter{General thoughts on scientific practices} \label{general}

Science is about discovering truths about the world around us and how we act in it. We can make mistakes or go down wrong paths, but we continually strive to improve our ideas and test our hypotheses so that we can ultimately increase our knowledge. The suggestions in this document will hopefully push you towards performing the process of discovery in the best way possible.

Most people don't set out to intentionally skew results or act for personal gain rather than scientific truth. Nevertheless, it's important to realize that many of our intuitions and common practices can lead to distorting the pattern or interpretation of results. Science is still a human an endeavor and will always be prone to mistakes and biases. The important question is how can we minimize bias/error and maximize the usefulness and truthiness of our results? This is not always easy, but will hopefully provide us with better---more accurate, reliable, and reproducible---results.

%%%%%%%%%%%%%%%%%%%%%%%%
\chapter{Making data and analyses publicly available} \label{open}

Sharing materials and data promotes more accurate science because (a) it facilitates checking of one's data and analysis by others, and (b) our internal checks on data organization and accuracy are typically better if we know the materials will be publicly available. So, even if no one ever looks at our data, the mere fact that we are sharing our data will improve the accuracy of our research. It may slow down the process, but this simply reflects a speed/accuracy tradeoff shifting towards more accuracy---which is a good thing for scientific discovery.

Making materials publicly available can save time down the road: most of us agree to make data available ``upon request'', but this can take many frustrating hours if the files have been moved or if the person responsible for a research project has moved on. Although organizing materials and data takes time now, it can save \textit{a lot} of time later.

When sharing data, it's best to share it on a third party site, rather than a personal or institutional website. Data posted online often becomes inaccessible due to websites disappearing, links changing, etc. Third party hosting sites are a safer bet for the long-term availability of your materials.

Many datasets can easily be shared using websites such as \href{https://github.com}{GitHub}, \href{https://figshare.com}{figshare}, or \href{https://datadryad.org/}{Dryad}. 

\vspace{1em} \noindent See also:

\begin{itemize}
\item Morey R et al.\ (2016) The Peer Reviewers' openness Initiative: Incentivizing open research practices through peer review. \journal{Royal Society Open Science} 3:150547. \doi{10.1098/rsos.150547}
\item Rouder JN (2016) The what, why, and how of born-open data. \journal{Behavior Research Methods} 48:1062--1069. \doi{10.3758/s13428-015-0630-z}
\item Vines, TH et al.\ (2014) The availability of research data declines rapidly with article age. \journal{Current Biology} 24, 94--97. \doi{10.1016/j.cub.2013.11.014}
\end{itemize}

%%%%%%%%%%%%%%%%%%%%%%%%
\chapter{Open-sourcing figures} \label{CCFigure}

Most journals will copyright the content of an article, including its figures, meaning that permission must be requested from the journal for any future use. In some (many?) cases this will be granted for free for academic purposes, but not always. Even if a paper isn't published in an open access journal, as the author you can typically license your figures before submitting to the journal. For example, if you use the popular \href{http://creativecommons.org/licenses/by/4.0/}{Creative Commons Attribution-By} license, the work is already licensed and so the journal does not own the copyright to the figure. Other authors (including you) are then free to easily re-use your figures (for example, in review papers), provided they provide attribution, which is generally what we want in science. 

%%%%%%%%%%%%%%%%%%%%%%%%
\chapter{Submitting preprints} \label{preprint}

Preprints generally refer to author-formatted manuscripts submitted to websites without pre-publication peer-review, which can be accessed by anyone. Popular preprint servers include \href{http://arxiv.org}{arXiv},  \href{http://biorxiv.org}{bioRxiv}, and \href{http://peerj.com}{PeerJ}.

The argument for preprints can be summarized in the following way: traditional scientific publishing takes too long, and thus delays the publicizing of research (which slows down science, and unfairly prevents authors from staking a claim to their ideas). In addition, research is frequently published in journals that charge for access, restricting who can benefit from the science. Finally, publishing a preprint may allow authors to get faster feedback on their work, resulting in a higher quality final product.

Some have argued against preprints. Like many things that go against traditional models preprints may be riskier for early stage scientists and trainees, particularly those not coming out of ``superstar'' labs. And, if comments are public, then feedback can also be public, which may be intimidating. Preprints may not be appropriate for every manuscript, but should at least be considered. 

In our lab, there is frequently no compelling reason to avoid submitting a preprint, apart from the extra work involved and the fact that it is scary.

\vspace{1em} \noindent See also:

\begin{itemize}
\item Bhalla N (2016) Has the time come for preprints in biology? \journal{Molecular Biology of the Cell} 27:1185--1187. \url{https://doi.org/10.1091/mbc.e16-02-0123}
\item Niko Kriegeskorte's ``The selfish scientist's guide to preprint posting'' \url{http://nikokriegeskorte.org/2016/03/13/the-selfish-scientists-guide-to-preprint-posting/}
\end{itemize}

%%%%%%%%%%%%%%%%%%%%%%%%
\chapter{Publishing in open access journals} \label{openaccess}

There are different flavors of open access, but the basic idea is that once a paper is published, anyone can read it. Open access contrasts with traditional publishing models, in which once a paper is published only readers who have paid for access to the journal or have institutional access can read it.

One challenge to open access is the publication fee, which is sometimes higher than for a traditional journal. A second challenge is that some open access journals are perceived to be of lower quality than their traditional counterparts. So, the theoretical and practical advantages of an open access publication must sometimes be tempered by these concerns.

\vspace{1em} \noindent See also:

\begin{itemize}
\item Martin Eve's ``Open access in a time of illness'' \newline \url{https://www.martineve.com/2016/04/07/open-access-in-a-time-of-illness/}
\item Tennant JP, Waldner F, Jacques DC et al.\ (2016) The academic, economic and societal impacts of Open Access: an evidence-based review. \journal{F1000Research} 5:632.\newline \doi{10.12688/f1000research.8460.1}
\end{itemize}


%%%%%%%%%%%%%%%%%%%%%%%%
\chapter{Explaining all conditions, tests, and data exclusions} \label{explainAllTests}

In a null hypothesis significance testing (NHST) framework, the more tests conducted, the better the chance of finding a significant effect even when one doesn't exist. For example, with a typical cutoff of p $<$ .05, if I conduct 20 tests, I am likely to find one ``significant'' difference by chance (when no true difference exists). Thus, it's important to be clear about how many tests were conducted, whether any data were excluded, and so on. (As researchers this should also mean that we are clear in our own minds when we are indeed testing a hypothesis vs.\ conducting exploratory research.)

Simmons et al.\ (2011) demonstrated this with example experiments and simulations with a large number of ``experimenter degrees of freedom'' to demonstrate that with enough tests it is possible to find ``significant'' results, and even to write these up in a convincing way. Explanations can be especially convincing when all of the tests aren't described, but only the ones leading a significant result. To avoid these situations, Simmons et al.\ (2012) suggest a ``21 word solution'' that makes it clear to readers whether all tests were (or were not) explained in the text:

\begin{quote}
We report how we determined our sample size, all data exclusions (if any), all manipulations, and all measures in the study.
\end{quote}

\vspace{1em} \noindent See also:

\begin{itemize}
\item Video from Neuroskeptic on p-hacking: \url{https://www.youtube.com/watch?v=A0vEGuOMTyA} 
\item Dorothy Bishop's take, featuring ``The amazing Significo'' \url{http://deevybee.blogspot.co.uk/2016/01/the-amazing-significo-why-researchers.html}
\item Carp J (2012) On the plurality of (methodological) worlds: Estimating the analytic flexibility of fMRI experiments. \journal{Frontiers in Neuroscience} 6:1-13. \doi{10.3389/fnins.2012.00149}
\item Simmons JP, Nelson LD, Simonsohn U (2011) False-positive psychology: Undisclosed flexibility in data collection and analysis allows presenting anything as significant. \journal{Psychological Science} 22:1359--1366. \doi{10.1177/0956797611417632}
\item Simmons JP,  Nelson LD, Simonsohn U (2012) A 21 Word Solution. \doi{10.2139/ssrn.2160588}
\end{itemize}


%%%%%%%%%%%%%%%%%%%%%%%%
\chapter{Collecting an informative data set} \label{n}

As discussed in the section on effect sizes, in many cases we don't just want to know whether a non-zero effect exists; we'd like to know \textit{how large} the effect is. For example, if I am studying the effect of hearing loss on speech intelligibility, I'd like to know \textit{how much} of a change in behavior I might see for a 5 dB change in hearing: it matters whether there is a 25\% drop in accuracy or a 1\% drop in accuracy.

Unfortunately, many power analyses (and more prevalent rules of thumb) are structured around finding a significant result, and not accurately estimating parameters. The number of participants needed to accurately estimate parameters, of course, varies, but estimates are typically in the range of hundreds of participants, rather than tens of participants. So---based on my experience and reading---generating an informative, accurate, and hopefully reproducible dataset typically requires many more participants than is typical for the field.

\vspace{1em} \noindent See also:

\begin{itemize}
\item Button KS, Ioannidis JPA, Mokrysz C, Nosek BA, Flint J, Robinson ESJ, Munaf� MR (2013) Power failure: why small sample size undermines the reliability of neuroscience. \journal{Nature Reviews Neuroscience} 14:365-376. \doi{10.1038/nrn3475}

\item Ioannidis JPA (2005) Why most published research findings are false. \journal{PLoS Medicine} 2:e124. \doi{10.1371/journal.pmed.0020124}

\item Kelley K, Maxwell SE (2003) Sample size for multiple regression: Obtaining regression coefficients that are accurate, not simply significant. \journal{Psychological Methods} 8:305-321. \doi{10.1037/1082-989X.8.3.305}

\item Lakens D, Evers ERK (2014) Sailing from the seas of chaos into the corridor of stability: Practical recommendations to increase the informational value of studies. \journal{Perspectives on Psychological Science} 9:278-292. \doi{10.1177/1745691614528520}

\item Maxwell SE (2000) Sample size and multiple regression analysis. \journal{Psychological Methods} 5:434-458. \doi{10.1037/1082-989X.5.4.434}

\item Sch\"{o}nbrodt F, Perugini M (2013) At what sample size do correlations stabilize? \journal{Journal of Research in Personality} 47:609-612. \doi{10.1016/j.jrp.2013.05.009}

\end{itemize}


%%%%%%%%%%%%%%%%%%%%%%%%
\chapter{Plotting data informatively} \label{plotting}

Bar graphs have been the standard in many years for showing means and effect sizes, but do a poor job of showing the distribution of the data: data with a myriad of underlying distributions can give rise to identical means (and standard deviations). With modern graphing software (e.g., R or Matlab) it's trivial to produce more informative plots, showing either every data point, or at least a fuller picture of the distribution (for example, in a \href{https://peerj.com/preprints/27137v1/}{raincloud plot}). Showing all the data gives readers a fuller set of the truth of the dataset (including the validity of any underlying assumptions that play in to the statistical analyses).


\vspace{1em} \noindent See also:

\begin{itemize}

\item Drummond GB, Vowler SL (2011) Show the data, don't conceal them. \journal{Journal of Physiology} 589.8:1861-1863. \doi{10.1113/jphysiol.2011.205062}

\end{itemize}


%%%%%%%%%%%%%%%%%%%%%%%%
\chapter{Potentially avoiding NHST and p values} \label{avoidNHST}

Writing in 1994, Cohen said:

\begin{quote}
Like many men my age, I mostly grouse. My harangue today is on testing for statistical significance, about which Bill Rozeboom (1960) wrote 33 years ago, ``The statistical folkways of a more primitive past continue to dominate the local scene'' (p. 417).

And today, they continue to continue. And we, as teachers, consultants, authors, and otherwise perpetrators of quantitative methods, are responsible for the ritualization of null hypothesis significance testing (NHST; I resisted the temptation to call it statistical hypothesis inference testing) to the point of meaninglessness and beyond. I argue herein that NHST has not only failed to support the advance of psychology as a science but also has seriously impeded it.
\end{quote}

\noindent More than 20 years later, the situation is largely unchanged. Difficulties with NHST include:

\begin{itemize}
\item The true null hypothesis is always false (i.e., means between two conditions will not be identical to an infinite precision)
\item A dichotomous estimate (``statistical significance'') rather than an interval
\item Misinterpretation of what p values reflect
\end{itemize}

\vspace{1em} \noindent See also:

\begin{itemize}

\item Geoff Cumming's ``Dance of the p values'' \url{https://www.youtube.com/watch?v=5OL1RqHrZQ8}
\item Cohen J (1994) The earth is round (p < .05). \journal{American Psychologist} 49:997--1003.\newline \doi{10.1037/0003-066X.49.12.997}
\item Loftus GR (1996) Psychology will be a much better science when we change the way we analyze data. \journal{Current Directions in Psychological Science} 5:161--171. \url{ http://www.jstor.org/stable/20182423}
\item Taleb NN (2016) The meta-distribution of standard p-values. \url{http://arxiv.org/abs/1603.07532}
\item Wagenmakers EJ, Morey RD, Lee MD (2016) Bayesian benefits for the pragmatic researcher. \url{https://webfiles.uci.edu/mdlee/BayesianBenefitsSubm.pdf}
\item Wagenmakers EJ, Verhagen AJ, Ly A (2016) How to quantify the evidence for the absence of a correlation. \journal{Behavior Research Methods} 48:413--426. \doi{10.3758/s13428-015-0593-0}
\item Wasserstein RL, Lazar NA (2016) The ASA's statement on p-values: Context, process, and purpose. \journal{The American Statistician}. 70:129--133. \doi{10.1080/00031305.2016.1154108}

\end{itemize}


%%%%%%%%%%%%%%%%%%%%%%%%
\chapter{Using confidence intervals} \label{CI}

The first point here is that any sort of interval or distribution suggests (quite correctly) that any number our statistical models produce is an \textit{estimate} rather than the ``truth''. You would react very differently if I said ``I will give you \$100'' compared to ``I will give you an amount between \$0 and \$210, or maybe I will even ask you to \textit{give} me \$10.''

Similarly, in a research study I may find that hearing loss is related to participants' memory for speech, with a standardized beta of 0.3. How much confidence do I have in this number of 0.3? It matters if the likely range is 0.28--0.32, or if it might be 0.0--0.6 (particularly because 0.0 would not be very interesting!).

So, thinking about distributions rather than point estimates is probably closer to what we actually care about.

Given that we care about distributions, there are many ways to report uncertainty around an estimate. The standard error (commonly used for error bars) gives an indication of the standard deviation of the population mean ( i.e., a distribution of sample means). But a point estimate, plus and minus 1 standard error, is only 68\% likely to contain the true mean.

Confidence intervals, on the other hand, should contain the population mean with a given accuracy (that is, the population mean should be within the 95\% confidence interval 95\% of the time). Thus, confidence intervals (CIs) are better suited for inference (that is, for drawing conclusions based on data). Reporting CIs instead of point estimates, and using CIs instead of standard errors, is a step in the right direction.

That being said, Morey et al.\ (2016) argue that most of us have an incorrect understanding of confidence intervals, which can have nonintuitive and/or unexpected properties. They suggest using Bayesian versions (credible intervals) instead. At the very least, informed by their simulations, it is worth checking assumptions about whether a confidence interval functions as we expect it to.

\vspace{1em} \noindent See also:

\begin{itemize}
\item Cumming G (2014) The new statistics: Why and how. \journal{Psychological Science} 25:7--29.\newline \doi{10.1177/0956797613504966}

\item Morey RD, Hoekstra R, Rouder JN, Lee MD, Wagenmakers E-J (2016) The fallacy of placing confidence in confidence intervals. \journal{Psychonomic Bulletin and Review} 23:103-123. \doi{10.3758/s13423-015-0947-8}


\end{itemize}



\end{document}